% Track F, stage 0 — candidate section for the report.
% Numbers from results/f01_stage0/summary.json at commit 3aa8a13 (run provenance in the file).
% Assets: report_candidates/tab_e11_pixel-path-verdict_{1,2}.tex (built by
% experiments/f01_tables.py). Do not edit the main LaTeX files from this branch.

\section{The pixel-path verdict without a second derivative}
\label{sec:pixel-path-stage0}

\paragraph{Question.}
When the flight path has to come from the pixels, the earlier pilot failed for two reasons
that are easy to conflate: it differentiated the recovered path twice, where noise grows as
frequency squared, and it needed a metric scale. This section asks whether five statistics
that avoid the second derivative can carry a verdict on the stored UZH-FPV DPVO paths, and
separates path error from depicted-lean error with an oracle ladder. It is pre-registered in
\texttt{configs/f01\_pixel\_path\_prereg.yaml} (commit \texttt{6b88d5a}), with one amendment
(\texttt{f01\_pixel\_path\_prereg\_amend1.yaml}, \texttt{17d21c6}) that fixes the path scale
and yaw offset on each clip's \emph{real} condition before its edits are scored: with a free
scale the tilt-removed condition is absorbed by $s\to0$ and the test cannot see it.

\paragraph{Setup.}
The path is Track C's DPVO run on the prepared UZH pinhole frames (150 frames per sequence,
15.07\,Hz, Kalibr intrinsics supplied; direction error 1.5--20.9$^\circ$ against truth). The
depicted lean is Track E's v3 fixed-crop GeoCalib output through the published Kalibr mount,
at the official schema-v2 0.37\,s window. The attitude source supplies the up direction
(truth quaternion, or the depicted lean inverted to a body-frame up) and a heading; a constant
yaw offset and a positive scale are fitted per clip, because yaw and scale are unconstrained in
this project.

The five statistics are: \textbf{V2 shooting} (integrate $a_h = g\,\mathrm{Rz}(\psi+\text{yaw})\,l_d
- k_d v_h$ twice and least-squares the observed path \emph{undifferentiated}; statistic is the
unexplained fraction $\rho$), \textbf{V1 preintegration} (one derivative of the path, integrate
the same ODE once, statistic $\rho_v$), \textbf{V3 band-limited coherence}
($1-|\rho_{\mathbb C}|$ between the depicted lean and the path's horizontal acceleration inside
$[0.05, f_{95}]$\,Hz), \textbf{V4 velocity-drag} ($l \approx k_d v/g$, one derivative), and
\textbf{V5 physics scale} (the scale the shooting fit recovers, checked for identifiability
and against a depth prior). V2 is primary because it uses the path undifferentiated.
Thresholds are leave-one-sequence-out conformal 95th percentiles on real clips; edits are
scored against the threshold of their held-out sequence, so no edited clip helps set its own
bar. The population is the eight real clips in four indoor-forward sequences that Track E
finds testable \emph{and} that the stored path covers in full; intervals are
sequence-clustered (5{,}000 replicates) and descriptive at four clusters.

\begin{table}[t]
\centering
\caption{Pixel-path verdict, estimated attitude and the stored DPVO path, 0.37\,s window.
Rates are flags and are printed as rate ($k/n$); $n=8$ real clips in four indoor-forward
sequences. Thresholds are leave-one-sequence-out conformal 95th percentiles on the real
clips. V4 is untestable here: its pre-registered drag-dominated condition
($\mathrm{median}\,|a_h|/(k_d|v_h|)<0.5$) holds on no clip (observed 0.73--9.35).}
\label{tab:pixel-path-primary}
\footnotesize
\begin{tabular}{lrrrrrrr}
\toprule
variant & real & null & yaw & removal & mirror & $0.5\times$ & $1.5\times$ \\
\midrule
V1 preintegration & 0.125 (1/8) & 0.125 (1/8) & 0.250 (2/8) & 0.125 (1/8) & 0.625 (5/8) & 0.250 (2/8) & 0.750 (6/8) \\
V2 shooting (primary) & 0.125 (1/8) & 0.125 (1/8) & 0.250 (2/8) & 0.750 (6/8) & 0.500 (4/8) & 0.375 (3/8) & 0.750 (6/8) \\
V3 band coherence & 0.250 (2/8) & 0.250 (2/8) & 0.125 (1/8) & 0.125 (1/8) & 0.000 (0/8) & 0.250 (2/8) & 0.125 (1/8) \\
V4 velocity-drag & -- & -- & -- & -- & -- & -- & -- \\
V5 physics scale & 0.000 (0/8) & 0.000 (0/8) & 0.000 (0/8) & 0.250 (2/8) & 0.000 (0/8) & 0.000 (0/8) & 0.000 (0/8) \\
\bottomrule
\end{tabular}

\end{table}

\begin{table}[t]
\centering
\caption{The oracle ladder for the shooting statistic: the same fit with the truth path in
place of the pixel path, and with the truth attitude in place of the depicted lean. Swap the
path and removal/mirror detection returns; swap the attitude and every edit becomes invisible
(the attitude arm then never reads the pixels). Rates are rate ($k/n$) over the same eight
clips.}
\label{tab:pixel-path-ladder}
\footnotesize
\begin{tabular}{llrrr}
\toprule
path & attitude & real & removal & mirror \\
\midrule
pixel & estimated attitude & 0.125 (1/8) & 0.750 (6/8) & 0.500 (4/8) \\
pixel & oracle attitude & 0.125 (1/8) & 0.125 (1/8) & 0.125 (1/8) \\
oracle & estimated attitude & 0.125 (1/8) & 1.000 (8/8) & 0.875 (7/8) \\
oracle & oracle attitude & 0.125 (1/8) & 0.125 (1/8) & 0.125 (1/8) \\
\bottomrule
\end{tabular}

\end{table}

\paragraph{Result.}
Table~\ref{tab:pixel-path-primary} gives the gate. With estimated attitude and the pixel path
the primary statistic flags 1/8 real clips (0.125), 6/8 tilt removals (0.750) and 4/8 mirrored
rolls (0.500); the null edit sits at 1/8 and the yaw control at 2/8. The false-reject budget
is met, the detection bar of 0.90 is not, and no secondary variant does better: preintegration
detects 1/8 removals and 5/8 mirrors, the band-limited coherence 1/8 and 0/8. The
\textbf{decision is fail}.

Table~\ref{tab:pixel-path-ladder} says why, and it is the useful part. With the truth path in
place of the pixel path the same code reaches 8/8 removal and 8/8 mirror under preintegration
and 8/8 and 7/8 under shooting, at a false-reject rate of 1--2/8. Swapping the attitude instead
leaves every condition identical, by construction, because the oracle-attitude arm never reads
the pixels. The statistic is therefore sound and the path is the binding error. Rank AUCs on
the same eight clips agree: shooting 0.95 for removal and 0.89 for mirror on the pixel path
against 1.00 and 0.98 on the truth path; preintegration 0.75/0.89 against 1.00/0.98. The gap
between those AUCs and the 0.75/0.50 rates in the table is the conformal threshold itself:
with four clusters a 95th percentile over six calibration clips is nearly their maximum, so a
statistic that separates well still misses the 0.90 detection bar. Lowering the level would
trade the (met) false-reject budget for detection; that was not pre-registered and is not done.

Two mechanical findings qualify the failure. First, detection collapses with the number of
derivatives of the observed path: no derivative (shooting) 0.75/0.50, one (preintegration)
0.125/0.625, two inside the band (coherence) 0.125/0.00. Second, the mirror rate is not ordered
by roll excitation in this sample (detected clips have roll spread 0.019--0.158, undetected
0.058--0.293), so the sample cannot support a claim that banking explains it; the per-clip
table is in the report. V5's ruler is identifiable on all eight real clips and its depth
consistency band holds, but it flags only 2/8 removals as unidentifiable and no scale edit,
because the pre-registered band is deliberately loose; the scale it recovers is 0.44--0.84
of the Sim(3) truth scale (median 0.59), a systematic bias the data here cannot attribute
between estimator attenuation and the drag/scale coupling.

\paragraph{Scope and caveats.}
The eight clips carry the gate; the exploratory set of fourteen further real clips (outdoor
rigs, truncated windows, the clip Track E scores but the path only partly covers) does not
separate at the pooled threshold either, with four testable real clips in its denominators.
The real path is reused for every edit condition: the v3 edits are pixel-space rotations, a
rotation about the optical centre leaves the camera's translation in its own frame unchanged,
and the one sequence where DPVO was run on edited frames (indoor\_forward\_5) reproduces the
real path to 4.98 units of rigid ATE on a 105.6-unit path with a 0.81$^\circ$ median step
direction difference. That is one sequence, and the estimator's sensitivity to the re-render
is otherwise unmeasured; it needs DPVO on the edited frames. Finally, with four clusters the
sequence-clustered intervals are descriptive, and one clip (indoor\_forward\_9\#01) is a false
reject whose edits are all flagged, so the unconditional detection rate is slightly flattering:
conditional on the real clip passing, shooting detects 5/7 removals and 3/7 mirrors.
